\input texsis
\paper

\overfullrule=0pt
\def\frac#1#2{#1\over #2}
%
\def\bmatrix#1{\left[ \matrix{#1} \right]}
\def\be{{\beta}}
\def\toone{{t+1}}

%\def\frac#1#2{#1\over #2}
%\magnification=1200
\overfullrule=0pt

\line{PhD Macroeconomics \hfil NYU}
\line{December 11,  2020 \hfil Thomas Sargent}
%
% \noindent Macroeconomics Q1 \hfill NYUAD \\
% Fall, 2019 \hfill Thomas Sargent
% \vspace{.3cm}

\centerline{\bf Final Examination}
%\centerline{\bf Homework 3}


\medskip

\noindent{\bf Instructions:}  This is a closed book  exam.  100 points are possible. Points attached to each problem are stated
at the beginning of each problem.

\medskip



\noindent{\bf 1. 25 points \quad} Where $c_t$ is consumption at date $t$ and $b_t$ is one-period debt owed at time $t$,
a consumer chooses a consumption, borrowing plan $\{c_t, b_{t+1}\}_{t=0}^T$
to maximize
$$ \sum_{t=0}^T\beta^t \ln(c_t),  \ \  \beta \in (0,1) $$
subject to the budget constraints
$$  c_t + b_t \leq y_t + \beta b_{t+1} , \ \ t = 0, \ldots, T $$
where $y_t = \lambda^t , \ \lambda \in (0,1)$, initial one-period risk-free debt $b_0 = 0$, and we impose
 $b_{T+1} \leq 0$.   


\medskip
\noindent{\bf a.} Verify that it is optimal for the consumer to set $b_{T+1} =0$.

\medskip
\noindent{\bf b.} Argue that  because the consumer can afford to set $c_t >0$ for all $t \in [0, 1, \ldots, T]$,
it follows that he will choose a  debt sequence that  satisfies 
$$ \eqalign{ b_T & \leq y_T \cr
                    b_{T-1} & \leq y_{T-1} + \beta y_T \cr
                    \vdots & \quad \quad \quad \vdots \cr
                    b_t & \leq y_t + \beta y_{t+1} + \cdots + \beta^{T-t} y_T }$$
                    
\medskip
\noindent{\bf c.} Argue that it is not necessary to impose so-called ``natural debt limits'' to make the consumer's
problem well posed.

\medskip
\noindent{\bf d.} Compute the consumer's optimal consumption plan.                                   


\bigskip

\noindent {\bf 2. 25 points \quad} Consider an infinite-horizon pure exchange economy that exists at $t=0, 1, \ldots $.
There are unit measures of  two types of consumers.
A type 1 consumer receives an endowment $y^1(s_t) = y(s_t) >0 $ of a nonstorable consumption good at times $t\geq T+1$
in Markov state $s_t \in {\bf S}$ and zero endowment for all $ t \leq T$.  A type 2 consumer receives an endowment $y^2(s_t) = y(s_t) >0$ of a nonstorable consumption good at times $t =0, \ldots, T$
in Markov state $s_t \in {\bf S}$ and zero endowment for all $t \geq T+1$.  Notice that $y^1(s_t) + y^2(s_t) = y(s_t)$ for all $t \geq 0$.

Complete markets in time- and  history-contingent consumption goods
meet at time $0$. A consumer of type 1  ranks consumption plans according to 
$$ \sum_{t=0}^T \sum_{s^t} \beta^t u(c_t^1(s^t)) \pi_t(s^t)$$
where $\pi_t(s^t)$ is a joint probability function over $s^t$ induced by the Markov chain $(P, \pi_0)$ for $s_t \in 
{\bf S}$.
A consumer of type 2 ranks consumption plans according to 
$$ \sum_{t=T+1}^\infty \sum_{s^t} \beta^t u(c_t^2(s^t)) \pi_t(s^t) .$$
%where $\pi_t(s^t)$ is a joint probability function over $s^t$ induced by the Markov chain $(P, \pi_0)$ on $s_t$.
A type $i$ consumer's budget constraint is
$$ \sum_{t=0}^\infty \sum_{s^t} q_t^0(s^t) c^i_t(s^t) \leq  \sum_{t=0}^\infty \sum_{s^t} q_t^0(s^t) y^i(s_t) . $$
 
\medskip
\noindent{\bf a.} Please define a competitive equilibrium with all trades occurring at time $0$.
\medskip
\noindent{\bf b.} Please compute a competitive equilibrium with time $0$ trades, being careful to compute
all objects that comprise a competitive equilibrium.
\medskip
\noindent{\bf c.} Please define an equilibrium with sequential trades of a complete set of one-period Arrow  securities.
\medskip
\noindent{\bf d.} Please compute an equilibrium with sequential trading of a complete set of one-period Arrow securities.

\bigskip

\noindent {\bf 3. 25 points \quad} An economy consists of overlapping generations of two-period live people. At dates $t \geq 1$ there are born
$N_t >0$ identical people who are endowed with $w_1 >0$ units of a nonstorable consumption good when they
are young and $w_2 >0 $ units of a nonstorable consumption good when they are old.  In addition, at time $1$ there
are $N_0> 0$ old people each of whom is endowed with ${M \over N_0}$ units of an intrinsically worthless unbacked (fiat) currency called
dollars (it is not in anyone's utility function and is useless for production).  
Population evolves as $N_{t+1} = n N_t$ for $t \geq 0$ where $n >0$.  The consumption profile
of a person born at time $t$ is $c_t^t$ when young and $c_{t+1}^t$ when old; he/she orders profiles according to the
utility functional
$ \ln (c_t^t) + \ln (c^t_{t+1}) .$

\medskip
\noindent 
{\bf a.} Define a competitive equilibrium in which no value is attached to fiat currency, being careful to define all
of the objects that comprise a competitive equilibrium and telling who trades what with whom.

\medskip
\noindent {\bf b.} Please compute all competitive equilibria in which no value is attached to fiat currency.

\medskip
\noindent{\bf c.} Please define a competitive equilibrium in which value is attached to fiat currency, 
being careful to define all
of the objects that comprise a competitive equilibrium and telling who trades what with whom.

\medskip
\noindent{\bf d.} A ``friend'' makes the following guess.  If $n w_1 > w_2$ there exist competitive equilibria
with valued fiat currency, while if $n w_1 < w_2$ there do not.  Please either confirm or dispose of this guess.
If the guess is true, please define and compute a {\it stationary\/} equilibrium with valued fiat currency and give formulas
for the equilibrium price level and rate of return on fiat currency.

\medskip
\noindent{\bf e.} Please argue that whenever a stationary equilibrium with valued fiat currency exists, then there
also exists a continuum of equilibria with valued fiat currency, in each of which the rate of return on currency
varies over time.  Please provide explicit formulas for an equilibrium  price level as a function of time.  
{\bf Hint:} It might help to deploy a linear difference equation.

\medskip
\noindent{\bf f.} Assuming that a nonstationary equilibrium with valued fiat currency exists, please compute the
the limiting value of the rate of return on currency as $t \rightarrow + \infty$. 


\medskip

\noindent {\bf 4. 25 points \quad} An economy consists of overlapping generations of two-period lived people.  Let $T > 2$ be a finite 
integer. At dates $ 1 \leq t \leq T$ there are born
$N >0 $ identical people who are endowed with $w_1 >0$ units of a nonstorable consumption good when they
are young and $w_2>0 $ units of a nonstorable consumption good when they are old.  
At dates $t \geq T+1$ there are born
$N$ identical people who are endowed with $\bar w_1 >0$ units of a nonstorable consumption good when they
are young and $\bar w_2>0$ units of a nonstorable consumption good when they are old. In addition, at time $1$ there
are $N$ old people each of whom is endowed with ${M \over N}$ units of an intrinsically worthless unbacked (fiat) currency called
dollars (it is not in anyone's utility function and is useless for production). The consumption profile
of a person born at time $t$ is $c_t^t$ when young and $c_{t+1}^t$ when old; he/she orders profiles according to the
utility functional
$ \ln (c_t^t) + \ln (c^t_{t+1}) .$


\medskip
\noindent 
{\bf a.} Define a competitive equilibrium in which no value is attached to fiat currency, being careful to define all
of the objects that comprise a competitive equilibrium and telling who trades what with whom.

\medskip
\noindent {\bf b.} Please compute all competitive equilibria in which no value is attached to fiat currency.

\medskip
\noindent{\bf c.} Please define a competitive equilibrium in which value is attached to fiat currency, 
being careful to define all
of the objects that comprise a competitive equilibrium and telling who trades what with whom.

\medskip
\noindent {\bf d.} Please state conditions on $w_1, w_2, \bar w_1, \bar w_2$ for which there exist competitive
equilibria with valued fiat currency.

\medskip
\noindent {\bf e.} Under sufficient conditions for there to exist competitive equilibria with valued fiat currency,
please find explicit formulas for the equilibrium price level sequence in the equilibrium for which the value
of currency is highest for all $t \geq 1$. {\bf Hint:} It might help to deploy a linear difference equation.




\medskip
\bye

